\section*{अध्याय \ref{ch:g5:solids} --- घन और घनाभ}

\begin{solution}{exo:g5:solids:1}
जैसे: एक पासा (\omterm{def:g5:solids:def}{घन}); अनाज का डिब्बा (\omterm{def:g5:solids:def}{घनाभ}); एक टिन (\omterm{def:g5:solids:def}{बेलन}); एक
फुटबॉल (गोला); कुछ छतें या पुराने मक़बरे
(\omterm{def:g5:solids:def}{पिरामिड})।
\end{solution}

\begin{solution}{exo:g5:solids:2}
\omterm{def:g5:solids:def}{घनाभ} के $6$ \omterm{def:g5:solids:def}{फलक}, $12$ \omterm{def:g5:solids:def}{किनारे}, $8$ \omterm{def:g5:solids:def}{शीर्ष} होते हैं --- और
$6 + 8 = 12 + 2$। \checkmark
\end{solution}

\begin{solution}{exo:g5:solids:3}
छहों \omterm{def:g5:solids:def}{फलक} आमने-सामने के एक जैसे फलकों के तीन जोड़ों में आते हैं
(ऊपर--नीचे, आगे--पीछे, बाएँ--दाएँ)। \omterm{def:g5:solids:def}{घन} पर छहों \omterm{def:g5:solids:def}{फलक} एक जैसे
वर्ग होते हैं।
\end{solution}

\begin{solution}{exo:g5:solids:4}
क्रॉस-जाल: पंक्ति में $2$ cm के चार वर्ग, दूसरे वर्ग के ऊपर एक
और नीचे एक। चारों की अँगूठी मोड़ने पर बग़लें बनती हैं; बाक़ी दो
वर्ग ऊपर और नीचे बंद कर देते हैं।
\end{solution}

\begin{solution}{exo:g5:solids:5}
क्रॉस-जाल पर (पंक्ति में चार वर्ग, दूसरे के ऊपर एक, नीचे एक)
आमने-सामने के जोड़े हैं: पंक्ति के वर्ग $1$ और $3$; पंक्ति के
वर्ग $2$ और $4$; ऊपर वाला और नीचे वाला वर्ग। एक सही नामांकन:
पंक्ति $= 1, 2, 6, 5$; ऊपर $= 3$; नीचे $= 4$ --- तब
$1{+}6$, $2{+}5$, $3{+}4$ सबका योग $7$ है।
\end{solution}

\begin{solution}{exo:g5:solids:6}
\omterm{def:g5:solids:def}{घनाभ}: हर परत में $5 \times 2 = 10$, $3$ परतें: $30$ \omterm{def:g5:solids:def}{घन}। \omterm{def:g5:solids:def}{घन}:
$3 \times 3 \times 3 = 27$। L-ढेरी: $3 \times 2 \times 1 = 6$ और
$1 \times 2 \times 1 = 2$: $8$ \omterm{def:g5:solids:def}{घन}।
\end{solution}

\begin{solution}{exo:g5:solids:7}
$12$ की परतें: जैसे $4 \times 3$ या $6 \times 2$ (या
$12 \times 1$)। $3$ परतों के साथ हर हाल में \omterm{def:g5:solids:volume}{आयतन}
$12 \times 3 = 36$ \omterm{def:g5:solids:def}{घन} है।
\end{solution}

\begin{solution}{exo:g5:solids:8}
$27 = 3 \times 3 \times 3$: $3$ भुजा वाला \omterm{def:g5:solids:def}{घन}।
$64 = 4 \times 4 \times 4$: $4$ भुजा वाला \omterm{def:g5:solids:def}{घन}।
\end{solution}

\begin{solution}{exo:g5:solids:9}
\omterm{def:g5:solids:def}{घनाभ} में $10 \times 5 \times 4 = 200$ चीनी के \omterm{def:g5:solids:def}{घन} आते हैं। रोज़
$6$ के हिसाब से: $200 = 6 \times 33 + 2$, इसलिए डिब्बा $33$ पूरे
दिन चलता है (और $34$वें दिन की सुबह के लिए $2$ \omterm{def:g5:solids:def}{घन} बचते हैं)।
\end{solution}

\begin{solution}{exo:g5:solids:10}
$3$ फलकों पर रंग: कोनों वाले $8$ \omterm{def:g5:solids:def}{घन}। $2$ फलकों पर: $12$ \omterm{def:g5:solids:def}{किनारों}
के बीच वाले: $12$ \omterm{def:g5:solids:def}{घन}। $1$ \omterm{def:g5:solids:def}{फलक} पर: $6$ फलकों के केंद्र वाले:
$6$ \omterm{def:g5:solids:def}{घन}। किसी पर नहीं: बीच में छिपा इकलौता \omterm{def:g5:solids:def}{घन}: $1$।
जाँच: $8 + 12 + 6 + 1 = 27$। \checkmark
\end{solution}
